%2multibyte Version: 5.50.0.2960 CodePage: 1252

\documentclass{article}
%%%%%%%%%%%%%%%%%%%%%%%%%%%%%%%%%%%%%%%%%%%%%%%%%%%%%%%%%%%%%%%%%%%%%%%%%%%%%%%%%%%%%%%%%%%%%%%%%%%%%%%%%%%%%%%%%%%%%%%%%%%%%%%%%%%%%%%%%%%%%%%%%%%%%%%%%%%%%%%%%%%%%%%%%%%%%%%%%%%%%%%%%%%%%%%%%%%%%%%%%%%%%%%%%%%%%%%%%%%%%%%%%%%%%%%%%%%%%%%%%%%%%%%%%%%%
\usepackage{amsmath}

\setcounter{MaxMatrixCols}{10}
%TCIDATA{OutputFilter=LATEX.DLL}
%TCIDATA{Version=5.50.0.2960}
%TCIDATA{Codepage=1252}
%TCIDATA{<META NAME="SaveForMode" CONTENT="1">}
%TCIDATA{BibliographyScheme=Manual}
%TCIDATA{Created=Monday, October 01, 2001 13:27:51}
%TCIDATA{LastRevised=Thursday, October 16, 2025 12:41:48}
%TCIDATA{<META NAME="GraphicsSave" CONTENT="32">}
%TCIDATA{<META NAME="DocumentShell" CONTENT="General\SW\Blank - Standard LaTeX Article">}
%TCIDATA{Language=American English}
%TCIDATA{CSTFile=LaTeX article (bright).cst}

\newtheorem{theorem}{Theorem}
\newtheorem{acknowledgement}[theorem]{Acknowledgement}
\newtheorem{algorithm}[theorem]{Algorithm}
\newtheorem{axiom}[theorem]{Axiom}
\newtheorem{case}[theorem]{Case}
\newtheorem{claim}[theorem]{Claim}
\newtheorem{conclusion}[theorem]{Conclusion}
\newtheorem{condition}[theorem]{Condition}
\newtheorem{conjecture}[theorem]{Conjecture}
\newtheorem{corollary}[theorem]{Corollary}
\newtheorem{criterion}[theorem]{Criterion}
\newtheorem{definition}[theorem]{Definition}
\newtheorem{example}[theorem]{Example}
\newtheorem{exercise}[theorem]{Exercise}
\newtheorem{lemma}[theorem]{Lemma}
\newtheorem{notation}[theorem]{Notation}
\newtheorem{problem}[theorem]{Problem}
\newtheorem{proposition}[theorem]{Proposition}
\newtheorem{remark}[theorem]{Remark}
\newtheorem{solution}[theorem]{Solution}
\newtheorem{summary}[theorem]{Summary}
\newenvironment{proof}[1][Proof]{\textbf{#1.} }{\ \rule{0.5em}{0.5em}}
\input{tcilatex}
\begin{document}


\begin{center}
{\LARGE Assignment 6}
\end{center}

\noindent Blankenau\hspace*{\stretch{0.7500}}

\noindent Economics 905, Fall 2025\medskip

\noindent \noindent \textit{Instructions.} You may work in groups and may
compare answers prior to submission. Joint submissions with up to 4 names
are allowed and encouraged. All work should be legible and concise. A due
date will be announced in class.

\begin{enumerate}
\item Describe how each of the following affects the $\dot{k}=0$ and $\dot{c}%
=0$ curves in our phase diagram, and thus how they affect the
balanced-growth-path values of $c$ and $k$:

\begin{enumerate}
\item A rise in $\theta .$

\item A change in the depreciation rate from the value of 0 assumed so far
to some positive level. That is, capital now evolves according to $\overset{.%
}{k}\left( t\right) =f\left( k\left( t\right) \right) -c\left( t\right)
-\left( n+g+\delta \right) k\left( t\right) .$ Note in addition that now the
return on capital is $r\left( t\right) =f^{\prime }\left( k\left( t\right)
\right) -\delta \ $rather than $r\left( t\right) =f^{\prime }\left( k\left(
t\right) \right) .$ This should lead you to conclude that both the $\dot{k}%
=0 $ and $\dot{c}=0$ shift.
\end{enumerate}

\item Consider a standard Ramsey-Cass-Koopman's growth model where
investment income is subsidized at rate $\varsigma _{k}.$ In each period,
revenue for the subsidy is generated from a lump sum tax of $G\left(
t\right) $ per effective unit of labor. The representative household
maximizes 
\begin{equation*}
U=\dint\limits_{t=0}^{\infty }e^{-\left( \rho -n-\left( 1-\theta \right)
g\right) t}\frac{c\left( t\right) ^{1-\theta }}{1-\theta }dt
\end{equation*}%
$\theta >0,$ $\ \rho -n-\left( 1-\theta \right) g>0$ subject to%
\begin{equation*}
k\left( 0\right) +\dint\limits_{t=0}^{\infty }e^{-R\left( t\right)
}e^{\left( n+g\right) t}\left( w\left( t\right) -G\left( t\right) \right)
dt=\dint\limits_{t=0}^{\infty }e^{-R\left( t\right) }e^{\left( n+g\right)
t}c\left( t\right) dt
\end{equation*}%
and the standard transversality condition. The capital stock per effective
unit of labor evolves according to $\overset{.}{k}\left( t\right) =f\left(
k\left( t\right) \right) -c\left( t\right) -\left( n+g\right) k\left(
t\right) .$ All items are defined as in the text: $c\left( t\right) $ and $%
k\left( t\right) $ are consumption and capital per effective unit of labor, $%
n$ in the growth rate of the population, $g$ is the growth rate of the labor
augmenting technology, $\rho $ is the discount rate, $R\left( t\right)
=\dint\limits_{\tau =0}^{t}r\left( \tau \right) d\tau $ $\ $where $r\left(
\tau \right) $ is the period $\tau $ interest rate, $w\left( t\right) $ is
the period $t$ wage. Thus the real interest rate facing households is given
by $\left( 1+\varsigma _{k}\right) f^{\prime }\left( k\left( t\right)
\right) .$ The intensive form production function $f\left( k\left( t\right)
\right) $ has the standard properties.

\begin{enumerate}
\item Show your derivation of the optimal growth rate of consumption.

\item Draw the $\overset{.}{k}\left( t\right) =0$ and $\overset{.}{c}\left(
t\right) =0$ loci in the $c,\,k$ space (that is, draw the principal curves
of the phase diagram). Label everything properly. Label the steady state
level of capital $k^{\ast }.$

\item For this part only let $n=g=0$ and $f\left( k\left( t\right) \right)
=y\left( t\right) =\left( k\left( t\right) \right) ^{.5}$. Suppose that the
economy has been operating at $k\left( t\right) =\frac{K_{0}}{L_{0}A_{0}}%
=k_{0}.$ Suddenly $k\left( t\right) $ changes to $k=zk_{0}\ $where $z<1$
This might have happened because of destruction of some of the capital
stock, a one-time permanent increase in $L$ or a one-time permanent increase
in $A$. A policymaker says that it does not matter what caused this decrease
since all dynamics are the same, regardless of the cause of the decrease. Do
you agree with this? If so, explain why. If not, describe the differences
across the three cases, paying particular attention to $y\left( t\right) $
and per capita output. Consider these items both immediately and in the long
run.

\item Suppose the economy is again at $k=k^{\ast }.$ What are the immediate
and long run consequence on $c$ and $k$ of an increase in $\varsigma _{k}?$
Show this on a phase diagram.

\item A policymaker says that in this situation, the effect of an increase
in $\varsigma _{k}=0$ to $\varsigma _{k}>0$ on utility, $U$, is ambiguous
because consumption is affected differently in the short run and in the long
run. Do you agree with this statement? Why or why not?\pagebreak 
\end{enumerate}
\end{enumerate}

\begin{enumerate}
\item 

\begin{enumerate}
\item ...For all of this, be sure to draw the pictures. I only describe them
here and can answer questions in class. From 
\begin{eqnarray*}
\dot{k}\left( t\right)  &=&f\left( k\left( t\right) \right) -c\left(
t\right) -\left( n+g\right) k\left( t\right)  \\
\frac{\dot{c}\left( t\right) }{c\left( t\right) } &=&\frac{f^{\prime }\left(
k\left( t\right) \right) -\rho -\theta g}{\theta }
\end{eqnarray*}%
we know that the the curves are  
\begin{eqnarray*}
\dot{k}\left( t\right)  &:&c\left( t\right) =f\left( k\left( t\right)
\right) -\left( n+g\right) k\left( t\right)  \\
\dot{c}\left( t\right)  &:&f^{\prime }\left( k\left( t\right) \right) =\rho
+\theta g
\end{eqnarray*}%
The parameter $\theta $ appears only in $f^{\prime }\left( k\left( t\right)
\right) =\rho +\theta g$ so only this curve moves. With $\rho +\theta g$
larger, $f^{\prime }\left( k\left( t\right) \right) $ must be larger which
occurs at a smaller $k\left( t\right) $. Thus this curve shifts left. The
new intesection is at a smaller c and k so steady state values of these are
lower. However, the immediate effect of $c\left( t\right) $ is a jump up to
get on the new saddle path. $k(t)$ does not change immediately. Note, then
that your new saddle bath should be above the $\dot{c}\left( t\right) =0$
curve. 

\item Now 
\begin{eqnarray*}
\dot{k}\left( t\right)  &:&c\left( t\right) =f\left( k\left( t\right)
\right) -\left( n+g+\delta \right) k\left( t\right)  \\
\dot{c}\left( t\right)  &:&f^{\prime }\left( k\left( t\right) \right) =\rho
+\theta g+\delta 
\end{eqnarray*}
Since $\delta $ is in each, both curves move. We consider them one at a
time. Your $\dot{k}\left( t\right) =0$ curve should shifts down. In
particluar, it intersects the horizonatal axis as the same point on the
left, but at a smaller amount of capital to the right so the new curve is
everywhere below the original curve. From this effect alone, the curves
intersect at the same capital and less consumption. There is an immedicate
jump to the new level of consumption. Immediate effects are less
consumption, same capital. Long run effect are less consumption, same
capital. Now consider the other effect shift in isolation. With $\rho
+\theta g+\delta $ larger $f^{\prime }\left( k\right) $ must be larger. This
occurs with less $k.$ The curve shifts to the left. Imediately, cousumption
increases to get on the new saddle path. In the long run, capital falls and
consumption falls since the new intersection is at lower $c$ and $k$. So
immediately, we have downward pressure on consumption from the first shift
and upward from the second. We do not know the effect on immediate
consumption because we do not know which of these dominates. Immediately,
niether moves capital so capital does not change. In the long run we have
downward pressure on consumption from both effects so it definately falls.
In the long run we have no change in capital from the first effect and a
downward change from the other. So $k$ definately falls in the long run. 
\end{enumerate}

\item ...

\begin{enumerate}
\item except for $r_{t}=\left( 1+\varsigma _{k}\right) f^{\prime }\left(
k\left( t\right) \right) ,$ this is the same as in your notes. 

\item as in notes. 

\item We know how to handle the case of larger $A$ from your notes, just
generalizing from $1.1$ to $z^{-1}$, teh relevant rows are     
\begin{equation*}
\begin{tabular}{llcc}
Item & Prior to shock & Immediate & Long Run \\ 
$y$ & \multicolumn{1}{|l}{$y_{0}=k_{0}^{.5}=\left( \frac{K_{0}}{A_{0}L}%
\right) ^{.5}$} & \multicolumn{1}{|c}{$\left( \frac{K_{0}}{A_{1}L}\right)
^{.5}=\left( \frac{1}{z^{-1}}\right) ^{.5}k_{0}^{.5}=z^{0.5}y_{0}$ } & 
\multicolumn{1}{|c}{$y_{1}=y_{0}$} \\ 
$\frac{Y}{L}$ & \multicolumn{1}{|l}{$A_{0}y_{0}$} & \multicolumn{1}{|c}{$%
z^{0.5}A_{1}y_{0}=z^{-0.5}A_{0}y_{0}$} & \multicolumn{1}{|c}{$%
z^{-1}A_{0}y_{0}$}%
\end{tabular}%
\end{equation*}%
\begin{tabular}{l|l|c|c}
\hline
$K$ & $K_{0}$ & $K_{0}$ & $K_{1}=z^{-1}K_{0}$ \\ 
$L$ & $L$ & $L$ & $L$ \\ 
$A$ & $A_{0}$ & $A_{1}=z^{-1}A_{0}$ & $A_{1}=z^{-1}A_{0}$ \\ 
$k$ & $k_{0}=\frac{K_{0}}{A_{0}L}$ & $\frac{K_{0}}{A_{1}L}=\frac{1}{z^{-1}}%
\frac{K_{0}}{A_{0}L}=zk_{0}$ & $k_{1}=k_{0}$%
\end{tabular}%
If it is instead a change in $L.$%
\begin{equation*}
\begin{tabular}{llcc}
Item & Prior to shock & Immediate & Long Run \\ \hline
$K$ & \multicolumn{1}{|l}{$K_{0}$} & \multicolumn{1}{|c}{$K_{0}$} & 
\multicolumn{1}{|c}{$K_{1}=z^{-1}K_{0}$} \\ 
$L$ & \multicolumn{1}{|l}{$L_{0}$} & \multicolumn{1}{|c}{$z^{-1}L_{0}$} & 
\multicolumn{1}{|c}{$z^{-1}L_{0}$} \\ 
$A$ & \multicolumn{1}{|l}{$A$} & \multicolumn{1}{|c}{$A$} & 
\multicolumn{1}{|c}{$A$} \\ 
$k$ & \multicolumn{1}{|l}{$k_{0}=\frac{K_{0}}{AL_{0}}$} & 
\multicolumn{1}{|c}{$\frac{K_{0}}{AL_{1}}=\frac{1}{z^{-1}}\frac{K_{0}}{AL_{0}%
}=zk_{0}$} & \multicolumn{1}{|c}{$k_{1}=k_{0}$} \\ 
$y$ & \multicolumn{1}{|l}{$y_{0}=k_{0}^{.5}=\left( \frac{K_{0}}{A_{0}L}%
\right) ^{.5}$} & \multicolumn{1}{|c}{$\left( \frac{K_{0}}{AL_{1}}\right)
^{.5}=\left( \frac{1}{z^{-1}}\right) ^{.5}k_{0}^{.5}=z^{0.5}y_{0}$ } & 
\multicolumn{1}{|c}{$y_{1}=y_{0}$} \\ 
$\frac{Y}{L}$ & \multicolumn{1}{|l}{$Ay_{0}$} & \multicolumn{1}{|c}{$%
z^{0.5}Ay_{0}$} & \multicolumn{1}{|c}{$Ay_{0}$}%
\end{tabular}%
\end{equation*}%
If it is instead $K$ 
\begin{equation*}
\begin{tabular}{llcc}
Item & Prior to shock & Immediate & Long Run \\ \hline
$K$ & \multicolumn{1}{|l}{$K_{0}$} & \multicolumn{1}{|c}{$zK_{0}$} & 
\multicolumn{1}{|c}{$K_{1}=K_{0}$} \\ 
$L$ & \multicolumn{1}{|l}{$L$} & \multicolumn{1}{|c}{$L$} & 
\multicolumn{1}{|c}{$L$} \\ 
$A$ & \multicolumn{1}{|l}{$A$} & \multicolumn{1}{|c}{$A$} & 
\multicolumn{1}{|c}{$A$} \\ 
$k$ & \multicolumn{1}{|l}{$k_{0}=\frac{K_{0}}{AL_{0}}$} & 
\multicolumn{1}{|c}{$\frac{zK_{0}}{AL}=zk_{0}$} & \multicolumn{1}{|c}{$%
k_{1}=k_{0}$} \\ 
$y$ & \multicolumn{1}{|l}{$y_{0}=k_{0}^{.5}=\left( \frac{K_{0}}{A_{0}L}%
\right) ^{.5}$} & \multicolumn{1}{|c}{$\left( \frac{zK_{0}}{AL}\right)
^{.5}=z^{0.5}y_{0}$ } & \multicolumn{1}{|c}{$y_{1}=y_{0}$} \\ 
$\frac{Y}{L}$ & \multicolumn{1}{|l}{$Ay_{0}$} & \multicolumn{1}{|c}{$%
z^{0.5}Ay_{0}$} & \multicolumn{1}{|c}{$Ay_{0}$}%
\end{tabular}%
\end{equation*}

FYI. This is hard to keep track of on an exam so a graphical and intuitve
response in that situation will suffice. 
\end{enumerate}
\end{enumerate}

\end{document}
